% ============================================================================
% Fichier  : partie6/P06_C25_impact_quantique.tex
% Titre    : Impact du Quantique sur l'Investigation Numérique
% Auteur   : MINKA MI NGUIDJOÏ Thierry Emmanuel
% Version  : 1.0 -- Juin 2026
% Statut   : RÉDIGÉ -- réécriture partielle
% Source partiellement rejetée : 9_impact_quantique_investigation_numerique.tex
%            contenait des éléments authentiques (algorithmes de Shor/Grover,
%            standards NIST PQC réels) mélangés à du code Python avec des
%            APIs fictives (pqcrypto.sign.dilithium2 dans cette forme
%            n'existe pas), des spéculations non vérifiables présentées
%            comme acquises ("quantum seal pour evidence bags",
%            "authentication quantique inviolable", détection QRNG par
%            seuils arbitraires non sourcés). Les faits techniques réels
%            ont été conservés et vérifiés ; le code fictif et les
%            spéculations non qualifiées ont été retirés ou explicitement
%            requalifiés comme hypothèses de recherche, non comme pratique
%            établie.
% Positionnement : Dernier chapitre du cours. S'appuie sur Ch.7 (ISO 27037,
%                  hash SHA-256), Ch.8 (custody, hash conservés séparément),
%                  Ch.11 (eIDAS, signatures qualifiées) pour poser la
%                  question de la pérennité de ces preuves face à une
%                  rupture cryptographique future.
% ============================================================================

\chapter{Impact du Quantique sur l'Investigation Numérique}
\label{chap:quantique}

\epigraph{%
    « Quantum computing will change everything we know about digital
    security and forensics. Preparation isn't optional --- it's
    essential. »%
}{--- Whitfield Diffie}

\niveauChapitre{Expert}{%
    Notions de cryptographie symétrique et asymétrique requises.
    Chapitres~\ref{chap:custody} (hash, intégrité) et
    \ref{chap:leg-mond} (eIDAS, signatures qualifiées) indispensables~:
    ce chapitre interroge leur pérennité à long terme.%
}

\begin{boxobjectifs}
\begin{itemize}
    \item Comprendre, sans détail mathématique excessif, pourquoi
          les algorithmes de Shor et Grover menacent la cryptographie
          actuelle et selon quel calendrier réaliste.
    \item Identifier précisément quels mécanismes forensiques du
          cours reposent sur une cryptographie potentiellement
          vulnérable~: hash, signatures, horodatage.
    \item Comprendre la menace ``\textit{Harvest Now, Decrypt Later}''
          et son impact sur les preuves numériques actuelles et futures.
    \item Connaître les standards de cryptographie post-quantique
          (PQC) effectivement normalisés par le NIST.
    \item Évaluer de façon critique les affirmations sur l'imminence
          de la menace quantique, en distinguant le consensus
          scientifique des spéculations commerciales.
    \item Formuler des recommandations pragmatiques pour la
          pratique forensique camerounaise actuelle.
\end{itemize}
\end{boxobjectifs}

\bigskip

% ============================================================================
\section*{Pourquoi ce chapitre clôt le cours}
% ============================================================================

Tout au long de ce cours, le hash SHA-256 a été présenté comme la
garantie ultime de l'intégrité d'une preuve (Ch.~\ref{chap:custody})~;
la signature électronique qualifiée comme la garantie ultime de
l'authenticité d'un document (Ch.~\ref{chap:leg-mond})~. Ces garanties
reposent sur une hypothèse implicite~: qu'aucun adversaire ne dispose
de la puissance de calcul nécessaire pour les casser. L'informatique
quantique, si elle atteint la maturité technique annoncée, lève
précisément cette hypothèse.

Ce chapitre ne prétend pas résoudre cette question~--- elle reste,
en juin~2026, un sujet de recherche actif et débattu. Il vise à vous
donner les moyens de comprendre l'enjeu, d'évaluer les affirmations
que vous rencontrerez sur ce sujet avec un esprit critique, et
d'adopter dès aujourd'hui des pratiques raisonnables sans céder
à un alarmisme commercial.

\separateur

% ============================================================================
\section{La Menace Quantique~: ce que Disent Réellement Shor et Grover}
\label{sec:menace-quantique}
% ============================================================================

\subsection{L'algorithme de Shor}

\begin{boxdef}
\textbf{Algorithme de Shor}\index{Shor, algorithme de} (Peter Shor,
1994)~: algorithme quantique permettant de factoriser de grands
nombres entiers et de résoudre le problème du logarithme discret
en temps polynomial --- contre un temps exponentiel pour les
meilleurs algorithmes classiques connus. Sa portée pratique dépend
entièrement de la disponibilité d'un \textbf{ordinateur quantique
suffisamment puissant et stable} pour l'exécuter sur des clés de
taille réelle, ce qui n'est pas le cas en~2026.
\end{boxdef}

L'algorithme de Shor menace directement les cryptosystèmes à clé
publique fondés sur la difficulté de factorisation ou du logarithme
discret~:

\begin{tableaucours}{p{3.0cm}p{3.5cm}p{4.5cm}}{%
    \tableauHeader{Algorithme} &
    \tableauHeader{Fondement mathématique} &
    \tableauHeader{Usage forensique concerné (ce cours)}
}
\textbf{RSA} & Factorisation de grands entiers &
    Signatures électroniques (Ch.~\ref{chap:leg-mond}),
    chiffrement TLS (Ch.~\ref{chap:for-res}) \\
\textbf{ECC / ECDSA} & Logarithme discret sur courbes elliptiques &
    Signatures électroniques modernes, Bitcoin (Ch.~\ref{chap:for-cloud})~;
    plus efficace que RSA, donc largement adopté \\
\textbf{Diffie-Hellman} & Logarithme discret &
    Établissement de clés pour connexions chiffrées (VPN, TLS) \\
\end{tableaucours}
\vspace{-0.8em}
\begin{center}{\small\itshape Tableau~25.1 --- Cryptosystèmes vulnérables à Shor}\end{center}

\begin{boxalert}
\textbf{Sur les chiffres de qubits~: prudence méthodologique}

Des estimations publiées dans la littérature scientifique évoquent
des ordres de grandeur de plusieurs milliers de \textbf{qubits
logiques} (et non physiques) nécessaires pour casser RSA-2048 dans
des délais opérationnels. Ces estimations varient significativement
selon les hypothèses techniques retenues (taux de correction
d'erreur, architecture). \textbf{Aucun consensus scientifique
ferme} n'existe en~2026 sur un calendrier précis. Méfiez-vous de
toute affirmation présentant une date exacte (``RSA cassé en~2035'')
comme un fait établi~: il s'agit, au mieux, de projections
spéculatives sujettes à révision constante.
\end{boxalert}

\subsection{L'algorithme de Grover}

\begin{boxdef}
\textbf{Algorithme de Grover}\index{Grover, algorithme de}
(Lov Grover, 1996)~: algorithme quantique offrant une
\textbf{accélération quadratique} (et non exponentielle) pour
la recherche dans une base de données non structurée. Appliqué
au cassage de clés symétriques, il réduit effectivement la
sécurité d'une clé de $n$~bits à environ~$n/2$~bits
``équivalent classique''.
\end{boxdef}

\begin{boxCRO}
\textbf{Analyse CRO --- Pourquoi Grover est une menace bien moindre que Shor}

L'accélération de Grover est \textbf{quadratique}, pas exponentielle.
Concrètement~: une clé AES-256 (utilisée dans le chiffrement
BitLocker du Ch.~\ref{chap:acq}, le principe 3-2-1 du
Ch.~\ref{chap:PIU}) offrirait, face à un ordinateur quantique
exploitant Grover, une sécurité ``équivalente'' à environ
128~bits classiques --- ce qui reste largement considéré comme
sûr selon les standards actuels.

$\mathcal{R}$~: la menace de Grover sur AES-256 est gérable
simplement en doublant la longueur de clé (AES-128~$\to$~AES-256)~;
elle ne nécessite pas de nouvel algorithme.

$\mathcal{O}$/$\mathcal{C}$~: la véritable rupture concerne la
cryptographie asymétrique (Shor), pas symétrique (Grover)~--- une
distinction souvent floutée dans les discours non spécialisés.

\cro{$\sim$ non affecté}{$\sim$ AES-256 reste robuste}{$\sim$ non affecté}

\textbf{Conséquence pratique~:} pour le chiffrement symétrique
utilisé en forensique (BitLocker, VeraCrypt, Ch.~\ref{chap:PIU}),
l'usage déjà recommandé d'AES-256 plutôt qu'AES-128 constitue
une protection raisonnable face à Grover. La vraie urgence
concerne les signatures et le chiffrement asymétrique.
\end{boxCRO}

% ============================================================================
\section{``Harvest Now, Decrypt Later''~: Implication pour les Preuves}
\label{sec:harvest-now}
% ============================================================================

\subsection{Le mécanisme de la menace}

\begin{boxdef}
\textbf{Harvest Now, Decrypt Later}\index{Harvest Now, Decrypt Later}~:
stratégie consistant, pour un adversaire, à intercepter et stocker
aujourd'hui des communications chiffrées qu'il ne peut pas
déchiffrer actuellement, dans l'attente qu'un ordinateur quantique
suffisamment puissant devienne disponible pour les déchiffrer
\textit{a posteriori}.
\end{boxdef}

Cette stratégie est documentée comme une préoccupation sérieuse
des agences de sécurité nationale (notamment américaines et
britanniques, qui ont publiquement justifié leurs feuilles de
route de migration PQC par ce risque), indépendamment de
l'incertitude sur le calendrier exact de la menace.

\subsection{Implication directe pour la forensique}

\begin{boxCRO}
\textbf{Analyse CRO --- Harvest Now, Decrypt Later et preuves numériques}

Cette menace inverse une intuition habituelle du cours~: jusqu'ici,
le risque pour une preuve portait sur son \textbf{intégrité présente}
(a-t-elle été altérée~?). La menace HNDL porte sur sa
\textbf{confidentialité future}~: des données chiffrées aujourd'hui
de façon réputée sûre (communications interceptées légalement,
Ch.~\ref{chap:for-res}~; données saisies chiffrées, Ch.~\ref{chap:acq})
pourraient devenir lisibles dans~10 à~20~ans par un tiers ayant
conservé une copie.

$\mathcal{C}$~: pour des données à \textbf{durée de sensibilité
longue} (affaires d'État, données de santé du Ch.~\ref{chap:for-iot},
secrets industriels), cette menace est directement pertinente
\emph{dès aujourd'hui}, car le chiffrement appliqué maintenant
devra rester sûr pendant des décennies.

$\mathcal{R}$/$\mathcal{O}$~: pour la plupart des affaires
forensiques courantes traitées dans ce cours (fraude MoMo,
ransomware, BEC), l'horizon temporel pertinent
(quelques années jusqu'au jugement) rend cette menace moins
urgente que pour des données à très longue durée de sensibilité.

\cro{$\downarrow\downarrow$ pour données longue durée}{$\sim$}{$\sim$}
\end{boxCRO}

% ============================================================================
\section{Impact sur la Chaîne de Custody et l'Intégrité Long Terme}
\label{sec:impact-custody}
% ============================================================================

\subsection{Trois mécanismes du cours potentiellement concernés}

\begin{tableaucours}{p{3.5cm}p{4.0cm}p{4.5cm}}{%
    \tableauHeader{Mécanisme (chapitre)} &
    \tableauHeader{Fondement cryptographique} &
    \tableauHeader{Vulnérabilité quantique}
}
Hash SHA-256
    (Ch.~\ref{chap:custody}) & Fonction de hachage --- pas de
    clé publique/privée &
    \textbf{Non vulnérable à Shor.} Affaibli par Grover, mais
    SHA-256 offre une marge de sécurité suffisante
    (cf.~Analyse CRO ci-dessus) \\
Signature électronique
    qualifiée eIDAS
    (Ch.~\ref{chap:leg-mond}) & Typiquement RSA ou ECDSA &
    \textbf{Vulnérable à Shor} si et quand un ordinateur quantique
    suffisant existe \\
Horodatage qualifié
    (Ch.~\ref{chap:leg-mond}) & Repose sur une signature
    (RSA/ECDSA) de l'autorité d'horodatage &
    \textbf{Vulnérable indirectement}, via la signature
    sous-jacente \\
\end{tableaucours}
\vspace{-0.8em}
\begin{center}{\small\itshape Tableau~25.2 --- Mécanismes forensiques et vulnérabilité quantique}\end{center}

\begin{boiteRemarque}{Une nuance importante~: hash et signature ne sont
pas équivalents face au quantique}
Une confusion fréquente mérite d'être dissipée~: le \textbf{hash
SHA-256} qui garantit l'intégrité d'une image E01 (Ch.~\ref{chap:acq})
n'est \textbf{pas} menacé de la même façon qu'une \textbf{signature
RSA/ECDSA} qui authentifie un document. Les fonctions de hachage
ne reposent pas sur la difficulté de factorisation~--- elles
résistent structurellement mieux à l'algorithme de Shor. C'est
la cryptographie asymétrique (signatures, échange de clés) qui
constitue le point de rupture principal, pas le hachage utilisé
quotidiennement en forensique.
\end{boiteRemarque}

\subsection{Le problème du re-timestamping}

Si l'algorithme de signature sous-jacent à un horodatage qualifié
devient un jour cassable, l'horodatage perd sa valeur probatoire
de \textit{preuve d'antériorité}. La pratique de
\textbf{re-timestamping périodique} (réappliquer un nouvel
horodatage, avec un algorithme à jour, sur une preuve archivée
avant que l'algorithme original ne soit compromis) est une
recommandation déjà formulée par certaines autorités
d'archivage à long terme (notamment dans le domaine de
l'archivage légal électronique), indépendamment du calendrier
précis de la menace quantique.

% ============================================================================
\section{Cryptographie Post-Quantique~: Standards Réellement Normalisés}
\label{sec:pqc-standards}
% ============================================================================

\subsection{Les standards NIST FIPS}

Contrairement à une partie significative du discours sur le sujet,
le NIST a effectivement publié des \textbf{standards finalisés}
de cryptographie post-quantique (PQC), et non de simples
propositions~:

\begin{tableaucours}{p{3.2cm}p{2.8cm}p{4.5cm}}{%
    \tableauHeader{Algorithme} &
    \tableauHeader{Standard NIST} &
    \tableauHeader{Usage}
}
\textbf{CRYSTALS-Kyber} & FIPS~203
    (\textit{ML-KEM}) & Encapsulation de clés (échange de
    clés sécurisé) --- standard principal \\
\textbf{CRYSTALS-Dilithium} & FIPS~204
    (\textit{ML-DSA}) & Signatures numériques --- standard
    principal \\
\textbf{SPHINCS+} & FIPS~205
    (\textit{SLH-DSA}) & Signatures fondées sur le hachage
    (alternative, plus conservatrice mathématiquement) \\
\textbf{FALCON} & En cours de
    standardisation & Signatures compactes, fondées sur les
    réseaux euclidiens (NTRU) \\
\end{tableaucours}
\vspace{-0.8em}
\begin{center}{\small\itshape Tableau~25.3 --- Standards PQC du NIST}\end{center}

\begin{boiteRemarque}{Ce qui distingue un standard authentique d'une
affirmation marketing}
La publication d'un \textbf{FIPS} (\textit{Federal Information
Processing Standard}) par le NIST est un processus public,
documenté, avec plusieurs années de revue par la communauté
cryptographique internationale. Quand vous rencontrez une
affirmation sur un algorithme ``post-quantique'', vérifiez
systématiquement~: s'agit-il d'un FIPS publié (vérifiable sur
\texttt{csrc.nist.gov}), d'un candidat encore en évaluation, ou
d'une simple allégation commerciale~? Cette vérification est
une application directe de l'esprit critique developpé tout au
long de ce cours (Ch.~\ref{chap:meth-std}~: distinguer un standard
réel d'une affirmation non sourcée).
\end{boiteRemarque}

\subsection{Migration hybride~: la pratique recommandée actuelle}

La transition recommandée par les agences de cybersécurité
occidentales (NSA, ANSSI, NCSC-UK) n'est \textbf{pas} un
remplacement immédiat et total de la cryptographie classique,
mais une approche \textbf{hybride}~: combiner un algorithme
classique éprouvé (RSA ou ECC) avec un algorithme post-quantique
normalisé, de sorte que la sécurité globale repose sur la
robustesse du \emph{plus fort} des deux mécanismes.

\begin{boxoutils}{Principe de la migration hybride (vue conceptuelle)}
\begin{lstlisting}[language=bash_forensique]
# PRINCIPE (illustration conceptuelle, pas une implementation a
# copier telle quelle -- toujours utiliser des bibliotheques
# cryptographiques auditees et maintenues, jamais de code
# cryptographique fait maison)

# Etablissement de cle hybride classique + post-quantique :
# 1. Generer un secret partage via ECDH (classique, eprouve)
# 2. Generer un secret partage via Kyber/ML-KEM (post-quantique)
# 3. Combiner les deux secrets (ex: via une fonction de derivation
#    de cle type HKDF) pour obtenir la cle finale
#
# Securite resultante = securite du MEILLEUR des deux mecanismes
# Un attaquant doit casser LES DEUX pour compromettre la cle

# Implementations reelles et auditees (a la date de redaction) :
# - OpenSSL 3.2+ : support experimental ML-KEM
# - liboqs (Open Quantum Safe project, OQS) : bibliotheque
#   de reference pour l'experimentation PQC
# - Signal Protocol : a deploye PQXDH (hybride X25519 + Kyber)
#   en production des 2023 pour l'echange de cles
\end{lstlisting}
\end{boxoutils}

% ============================================================================
\section{``Quantum Forensics''~: Distinguer Recherche et Pratique Établie}
\label{sec:quantum-forensics-critique}
% ============================================================================

\begin{boxalert}
\textbf{Mise en garde sur les affirmations de ``forensique quantique''}

Certains discours évoquent des concepts comme l'``authentification
quantique inviolable'', le ``scellement quantique de preuves''
(\textit{quantum seal}) ou la détection automatisée de génération
quantique aléatoire par de simples tests statistiques de seuil.

\textbf{Évaluation critique, conforme à la méthode de ce cours
(distinction constatation/analyse, Ch.~\ref{chap:meth-std})~:}
\begin{itemize}
    \item Ces concepts relèvent, en~2026, de la \textbf{recherche
          exploratoire}, pas de pratiques forensiques établies
          ou normalisées~;
    \item Aucun standard international (ISO, NIST) ne définit
          actuellement de procédure de ``scellement quantique''
          opérationnelle pour la chaîne de custody~;
    \item La distinction entre un générateur de nombres aléatoires
          quantique (QRNG) et un générateur pseudo-aléatoire
          classique (PRNG) sophistiqué par de simples tests
          statistiques de surface est un sujet de recherche actif
          et non trivial~--- une affirmation présentant cela comme
          une procédure forensique simple et fiable doit être
          accueillie avec un scepticisme appuyé.
\end{itemize}
\textbf{Posture recommandée~:} appliquer ici exactement le même
principe qu'au Ch.~\ref{chap:rapport} pour toute affirmation
technique dans un rapport forensique --- citer la source, le
niveau de maturité réel (norme publiée~? article de recherche~?
communiqué commercial~?), et ne jamais présenter une hypothèse
de recherche comme une pratique établie.
\end{boxalert}

% ============================================================================
\section{Recommandations Pragmatiques pour la Pratique Actuelle}
\label{sec:recommandations-pratiques}
% ============================================================================

\subsection{Ce qui est raisonnable de faire dès aujourd'hui}

\begin{enumerate}
    \item \textbf{Continuer à utiliser SHA-256 (voire SHA-512)}
          pour le hachage d'intégrité (Ch.~\ref{chap:custody})~:
          ces fonctions ne sont pas la priorité de la menace
          quantique et restent le standard recommandé.
    \item \textbf{Privilégier AES-256 à AES-128} pour le chiffrement
          symétrique de longue conservation (sauvegardes
          forensiques, Ch.~\ref{chap:PIU})~: marge de sécurité
          suffisante face à Grover, sans coût opérationnel
          significatif.
    \item \textbf{Pour les preuves à très longue durée de
          sensibilité} (affaires d'État, données de santé,
          secrets industriels, Ch.~\ref{chap:for-iot})~: envisager
          un re-timestamping périodique et suivre l'évolution
          des recommandations PQC des autorités internationales
          de référence.
    \item \textbf{Documenter dans le rapport forensique}
          (Ch.~\ref{chap:rapport}) l'algorithme de signature
          utilisé pour tout horodatage qualifié, afin de faciliter
          une future ré-authentification si nécessaire.
    \item \textbf{Ne pas céder à l'urgence commerciale}~: pour
          la quasi-totalité des affaires traitées dans ce cours
          (durée de pertinence~: quelques années), la menace
          quantique n'impose aucun changement de pratique immédiat
          au-delà des bonnes pratiques déjà recommandées
          (SHA-256, AES-256).
\end{enumerate}

\begin{boxjuris}
\textbf{Application au cadre camerounais}

La \loicam{2024/017} (Ch.~\ref{chap:droit-cam}) impose, en son
art.~27, des ``mesures techniques et organisationnelles
appropriées'' pour la sécurité des données. Cette obligation
de proportionnalité s'applique également à l'anticipation
quantique~: un responsable de traitement gérant des données de
santé ou des secrets d'État doit raisonnablement intégrer cette
menace à son analyse de risque~; un commerce de détail traitant
des données clients courantes n'a pas la même urgence. La
proportionnalité, principe déjà rencontré au Ch.~\ref{chap:droit-cam},
s'applique pleinement à la planification de la transition PQC.
\end{boxjuris}

\separateur

\begin{boxresume}
\textbf{Ce qu'il faut retenir de ce chapitre~:}
\begin{itemize}
    \item \textbf{Shor menace la cryptographie asymétrique}
          (RSA, ECC, signatures, échange de clés)~--- rupture
          potentiellement totale, mais sans calendrier scientifique
          consensuel précis.

    \item \textbf{Grover affaiblit la cryptographie symétrique}
          de façon \textbf{quadratique seulement}~: AES-256 reste
          robuste face à cette menace (équivalent~128~bits classiques).
          Ne pas confondre l'ampleur des deux menaces.

    \item \textbf{Le hash SHA-256} qui garantit l'intégrité des
          preuves (Ch.~\ref{chap:custody}) n'est \textbf{pas}
          structurellement menacé comme une signature RSA/ECDSA~:
          distinction essentielle souvent confondue.

    \item \textbf{Harvest Now, Decrypt Later}~: pertinent dès
          aujourd'hui pour les données à \emph{longue} durée de
          sensibilité~; moins urgent pour la majorité des affaires
          forensiques courantes de ce cours.

    \item \textbf{Standards PQC réels~:} FIPS~203 (Kyber/ML-KEM),
          FIPS~204 (Dilithium/ML-DSA), FIPS~205 (SPHINCS+/SLH-DSA)
          --- vérifiables sur \texttt{csrc.nist.gov}. La migration
          recommandée est \textbf{hybride}, pas un remplacement brutal.

    \item \textbf{Posture critique obligatoire}~: distinguer
          systématiquement un standard publié d'une affirmation
          commerciale ou d'une hypothèse de recherche
          (``quantum seal'', ``authentification quantique
          inviolable'')~--- même méthode que la distinction
          constatation/analyse du Ch.~\ref{chap:rapport}.

    \item \textbf{Recommandations pragmatiques}~: SHA-256/512
          et AES-256 restent appropriés aujourd'hui~; re-timestamping
          périodique pour les preuves à très longue durée de
          sensibilité~; documentation systématique des algorithmes
          utilisés~; proportionnalité du risque (art.~27
          Loi~2024/017).
\end{itemize}
\end{boxresume}

\bigskip

\begin{boxexo}[title={\ding{45}\quad Exercice~25.1 --- Distinguer les menaces \niveaubase}]
Pour chacun des mécanismes suivants, rencontrés dans ce cours,
indiquez s'il est principalement menacé par Shor, par Grover,
par les deux, ou par aucun des deux, et justifiez~:
\begin{enumerate}[(a)]
    \item Le hash SHA-256 d'une image E01 (Ch.~\ref{chap:acq}).
    \item Une signature électronique qualifiée eIDAS sur un
          document européen (Ch.~\ref{chap:leg-mond}).
    \item Le chiffrement BitLocker (AES-256) d'un disque dur
          saisi (Ch.~\ref{chap:PIU}).
    \item L'établissement d'une connexion TLS lors d'une capture
          réseau (Ch.~\ref{chap:for-res}).
\end{enumerate}
\end{boxexo}

\begin{boxexo}[title={\ding{45}\quad Exercice~25.2 --- Esprit critique appliqué \niveaumoyen}]
Un fournisseur de solutions de sécurité vous présente la publicité
suivante~: ``Notre coffre-fort numérique offre une authentification
quantique inviolable grâce à notre technologie de scellement
quantique propriétaire, garantissant une protection absolue
contre toute menace présente et future.''
\begin{enumerate}[(a)]
    \item Appliquez la grille d'évaluation critique de la
          Section~\ref{sec:quantum-forensics-critique}~: quelles
          questions précises poseriez-vous à ce fournisseur~?
    \item Quels éléments de cette publicité relèvent du marketing
          plutôt que d'un standard vérifiable~?
    \item Si vous deviez recommander une solution de protection
          réelle pour des preuves à longue durée de sensibilité,
          que recommanderiez-vous à la place, en vous appuyant
          sur les standards effectivement normalisés
          (Section~\ref{sec:pqc-standards})~?
\end{enumerate}
\end{boxexo}

\begin{boxexo}[title={\ding{45}\quad Exercice~25.3 --- Politique d'archivage à long terme \niveauexpert}]
L'hôpital de l'affaire SANAGA (Ch.~\ref{chap:cas-int}) doit conserver
les dossiers patients chiffrés pendant~30~ans conformément à la
réglementation sur les données de santé. Vous êtes consulté pour
définir une politique de protection cryptographique à long terme.
\begin{enumerate}[(a)]
    \item Appliquez le principe de proportionnalité
          (\loicam{2024/017}, art.~27)~: en quoi ce cas diffère-t-il
          d'une PME traitant des données clients courantes en
          termes d'urgence quantique~?
    \item Proposez une stratégie de chiffrement combinant
          chiffrement symétrique (pour les données elles-mêmes)
          et signature/horodatage (pour leur intégrité et
          antériorité), en justifiant chaque choix d'algorithme
          par les éléments de ce chapitre.
    \item À quelle fréquence recommanderiez-vous un re-timestamping~?
          Sur quels critères objectifs baseriez-vous cette
          fréquence (plutôt qu'une date arbitraire)~?
    \item Rédigez un paragraphe destiné au rapport d'audit de
          sécurité de l'hôpital, distinguant clairement ce qui
          relève de mesures établies et vérifiables de ce qui
          relève d'une anticipation prudente face à une incertitude
          scientifique persistante.
\end{enumerate}
\end{boxexo}

\bigskip

\begin{boxinfo}
\textbf{Conclusion du cours.} Ce chapitre clôt le parcours entamé
au Prologue de ce cours~: de l'affaire MBONGO, où une simple erreur
procédurale a anéanti un dossier judiciaire, à la question de la
pérennité cryptographique des preuves face à une rupture
technologique encore incertaine. Entre ces deux extrêmes, vous
avez acquis~: la rigueur méthodologique du PIU, la maîtrise
technique de l'acquisition et de l'analyse multi-supports, la
compréhension du cadre juridique camerounais et international,
et la posture éthique et critique indispensable à la pratique
de l'investigation numérique.

\textit{Cette dernière compétence --- l'esprit critique face aux
affirmations non vérifiées, qu'elles proviennent d'un suspect,
d'un outil automatisé, ou d'un discours commercial sur une
technologie émergente --- est sans doute la plus durable de
toutes celles enseignées dans ce cours. Les outils changeront.
Les lois évolueront. La rigueur intellectuelle, elle, reste.}
\end{boxinfo}